%!TEX program = lualatex
% SPDX-License-Identifier: CC-BY-SA-4.0
% Requires: https://github.com/jemmybutton/fiziko (GPL-3.0)
% Set MPINPUTS to the directory containing fiziko.mp before running lualatex.
\documentclass[10pt]{article}

\usepackage{iftex}
\ifLuaTeX\else
  \PackageError{vintage-latex}{This example requires LuaLaTeX}
    {Run lualatex, not pdflatex or xelatex.}
\fi

\usepackage{fontspec}
\usepackage{unicode-math}
\usepackage{xcolor}
\usepackage{geometry}
\usepackage{multicol}
\usepackage[tracking=true,protrusion=false,expansion=false]{microtype}
\usepackage{luamplib}

\setmainfont{EB Garamond}[Numbers={OldStyle,Proportional}]
\setmathfont{Garamond-Math.otf}
\definecolor{ragpaper}{HTML}{F3E9D6}
\definecolor{printerink}{HTML}{25201B}
\pagecolor{ragpaper}
\color{printerink}
\geometry{a4paper,top=22mm,bottom=24mm,left=22mm,right=22mm}
\pagestyle{empty}
\setlength{\parindent}{1em}
\setlength{\columnsep}{8mm}
\setlength{\columnseprule}{.3pt}
\setlength{\emergencystretch}{.8em}

\mplibforcehmode
\everymplib{
  input fiziko.mp;
  randomseed := 3141;
  defineMinStrokeWidth(0.2pt);
  defineLightDirection(-1/8pi, 1/8pi);
  color inkcolor; inkcolor := (0.145, 0.125, 0.106);
  drawoptions(withcolor inkcolor);
  beginfig(1);
}
\everyendmplib{endfig;}

\newcommand{\articlehead}[1]{%
  \par\medskip
  \begin{center}\textls[60]{\textsc{#1}}\end{center}
  \vspace{-.3em}
}
\newcommand{\figcaption}[1]{\par\vspace{.3em}\centerline{\small\itshape #1}\par\smallskip}

\begin{document}

\begin{center}
  {\small\textls[120]{\textsc{Chapter the Third}}}\par
  \vspace{.3em}
  {\Large Elements of Mechanics}\par
  \vspace{.2em}
  {\itshape illustrated by figures of the simplest engines}\par
  \vspace{.6em}
  \rule{.5\linewidth}{.4pt}
\end{center}

\vspace{.6em}

\begin{multicols}{2}

\articlehead{Of the Cart}

\noindent A load placed upon wheels is moved by a force much less than its
weight, because the wheels turn instead of rubbing along the ground. The
force $F$ has then only to overcome the small friction at the axles, together
with the resistance of the ground to the rolling of the rims.

\begin{center}
\begin{mplibcode}
  numeric l, w, r, h;
  l := 3.4cm; w := 0.26cm; r := 0.6cm; h := 0.9cm;
  draw solidSurface((-0.6cm, 0) -- (l + 1.5cm, 0));
  draw woodBlock(l, w) shifted (0, r);
  draw wheel(r, 0) shifted (r, 1/2r);
  draw wheel(r, 0) shifted (l - r, 1/2r);
  draw weight.s(h) shifted (1/2l, r + w);
  drawarrow (l, r + w + 1/2h) -- (l + 1.2cm, r + w + 1/2h)
    withpen pencircle scaled 0.45pt;
  label.top(btex $F$ etex, (l + 0.6cm, r + w + 1/2h + 0.1cm));
\end{mplibcode}
\figcaption{Fig.\ 1. A loaded cart drawn along level ground.}
\end{center}

\articlehead{Of the Lever}

\noindent A bar resting upon a prop, and carrying weights at its two ends,
is at rest when the weights are to one another inversely as their distances
from the prop. Thus in Fig.\ 2 the greater weight, being nearer the prop,
balances the smaller placed farther from it; and the rule is
\[
  W_1 \, a = W_2 \, b .
\]

\begin{center}
\begin{mplibcode}
  numeric tilt, barLength, barThick;
  pair pivot, leftEnd, rightEnd;
  path prop, bar;
  tilt := -5;
  barLength := 5.2cm; barThick := 0.22cm;
  pivot := (0, 0.95cm);
  draw solidSurface((-2.6cm, 0) -- (3.2cm, 0));
  prop := (-0.5cm, 0) -- (0.5cm, 0) -- (0, 0.95cm) -- cycle;
  draw woodenThing(prop, 90);
  bar := unitsquare xscaled barLength yscaled barThick
         shifted (-2/5barLength, 0) rotated tilt shifted pivot;
  draw woodBlock(barLength, barThick)
       shifted (-2/5barLength, 0) rotated tilt shifted pivot;
  leftEnd := (-2/5barLength + 0.35cm, barThick) rotated tilt shifted pivot;
  rightEnd := (3/5barLength - 0.35cm, barThick) rotated tilt shifted pivot;
  draw weight.s(0.85cm) shifted leftEnd;
  draw weight.s(0.55cm) shifted rightEnd;
  label.top(btex $W_1$ etex, leftEnd shifted (0, 0.9cm));
  label.top(btex $W_2$ etex, rightEnd shifted (0, 0.6cm));
  drawdblarrow (xpart leftEnd, -0.45cm) -- (0, -0.45cm)
    withpen pencircle scaled 0.35pt;
  drawdblarrow (0, -0.45cm) -- (xpart rightEnd, -0.45cm)
    withpen pencircle scaled 0.35pt;
  label.bot(btex $a$ etex, (1/2xpart leftEnd, -0.5cm));
  label.bot(btex $b$ etex, (1/2xpart rightEnd, -0.5cm));
  draw (xpart leftEnd, -0.6cm) -- (xpart leftEnd, -0.3cm)
    withpen pencircle scaled 0.3pt;
  draw (0, -0.6cm) -- (0, -0.3cm) withpen pencircle scaled 0.3pt;
  draw (xpart rightEnd, -0.6cm) -- (xpart rightEnd, -0.3cm)
    withpen pencircle scaled 0.3pt;
\end{mplibcode}
\figcaption{Fig.\ 2. Two weights in balance upon a lever.}
\end{center}

\columnbreak

\articlehead{Of the Pendulum}

\noindent A heavy ball hung by a thread, and drawn aside through the angle
$\theta$, returns toward the lowest point and passes beyond it, so that the
motion is repeated. For small excursions the time of one swing depends upon
the length $l$ alone, and not upon the weight of the ball nor upon the width
of the arc; whence the pendulum serves to divide the day into equal parts.

\begin{center}
\begin{mplibcode}
  numeric l, theta, ballDiam;
  pair top, bob;
  l := 3.6cm; theta := 24; ballDiam := 0.8cm;
  top := (0, 0);
  bob := top + l * dir(-90 + theta);
  draw solidSurface((1.5cm, 0) -- (-1.5cm, 0));
  draw top -- (0, -l - 1/2ballDiam - 0.3cm)
    dashed evenly scaled 0.6 withpen pencircle scaled 0.28pt;
  draw top -- bob withpen pencircle scaled 0.4pt;
  draw sphere.c(ballDiam) shifted bob shifted (1/2ballDiam * dir(-90 + theta));
  markAngle((0, -l), top, bob)(btex $\theta$ etex);
  label.rt(btex $l$ etex, 1/2[top, bob] shifted (0.12cm, 0));
  draw (subpath (5.35, 6.65) of fullcircle scaled (2l + ballDiam + 1.1cm))
    shifted top dashed withdots scaled 0.6 withpen pencircle scaled 0.35pt;
\end{mplibcode}
\figcaption{Fig.\ 3. A pendulum drawn aside from the vertical.}
\end{center}

\articlehead{Of Weight and Measure}

\noindent The figures are drawn by the same MetaPost library that made the
plates of the preceding chapters, but here they are produced inside the page
itself: LuaTeX's built-in MetaPost interpreter renders each one at the place
where it stands in the column, and the letters upon them are set in the type
of the text. The wood grain, the hatching of the ground, the shading of the
ball, and the standing weights are stock parts of the library; the arrows,
dimension lines, and dotted arc are plain MetaPost added around them.

\smallskip
\noindent A page of this kind was once the work of an engraver and a
compositor in turn. Here the two are done together, and a change of a single
length in the source moves the weights, the arrows, and the letters at once.

\end{multicols}

\vfill
\begin{center}
  {\scriptsize\textls[110]{\textsc{To be continued in the next number}}}
\end{center}

\end{document}
